\documentclass[a4paper,12pt]{article}

\input{../../Preambulo.tex}
\newcommand{\assunto}{FUNÇÕES QUADRÁTICAS}
\newcommand{\atividade}{atividade} % prova ou lista ou as ou ad ou atividade
\newcommand{\NumQObj}{6}
\newcommand{\ValorProva}{2}
\pgfmathparse{\ValorProva/\NumQObj}
\edef\ValorQObj{\pgfmathresult}
%\newcommand{\ValorQObj}{0.3}
%\pgfmathparse{int(round(\NumQObj * \ValorQObj))}
%\edef\ValorProva{\pgfmathresult}
\newcommand{\tipo}{dupla} % dupla ou individual
\newcommand{\calculadora}{nao} % sim ou nao
\newcommand{\turma}{ELETROELETRÔNICA 1.\textordmasculine\ ANO}
\newcommand{\numProva}{A}
\newcommand{\professor}{Igor Martins Silva}
\newcommand{\bimestre}{3}
\newcommand{\ano}{2026}
\newcommand{\mesTexto}{agosto}
\newcommand{\mesNum}{08}
\newcommand{\dia}{27}
\newcommand{\dataTexto}{\dia\ de \mesTexto\ de \ano}
\newcommand{\dataNun}{\dia/\mesNum/\ano}

\newcommand{\titulo}{}

\ifdefstring{\atividade}{prova}
  {\renewcommand{\titulo}{Avaliação de Matemática}}
{
\ifdefstring{\atividade}{lista}
  {\renewcommand{\titulo}{Lista de exercícios de Matemática}}
{
\ifdefstring{\atividade}{as}
  {\renewcommand{\titulo}{Avaliação Somativa}}
{
\ifdefstring{\atividade}{ad}
  {\renewcommand{\titulo}{Avaliação Diagnóstica}}
  {\renewcommand{\titulo}{Atividade de Matemática}}
}}}

\newcommand{\emtipo}{%
  \ifdefstring{\tipo}{individual}{\tipo}{em \tipo}%
}

\edef\pathCapa{\currfiledir}

\newcommand{\subtitulobase}[3]{%
    % #1 = título da 1ª coluna
    % #2 = conteúdo da 1ª coluna
    % #3 = mostra número da prova? (vazio ou algo)
    
    \noindent
    \begin{minipage}{2.8cm}
        \includegraphics[width=2.8cm]{\pathCapa Logo_Cefet.png}
    \end{minipage}
    \hfill
    \begin{minipage}{\dimexpr\textwidth-6.2cm\relax}
        \begin{center}
        CEFET-MG -- Campus Contagem \\[3pt]
        \titulo\ -- \bimestre.\textordmasculine\ bimestre de \ano \\[3pt]
        \professor
    \end{center}
    \end{minipage}
    \hfill
    \begin{minipage}{3.2cm}
        \includegraphics[width=3.2cm]{\pathCapa Brasao_Brasil.png}
    \end{minipage}
    
    \begin{center}
        \vspace{15pt}
        \begin{tabular}{|C{210pt}|C{70pt}|C{210pt}|}
            \hline
            \textbf{#1} &
            \textbf{DATA} &
            \textbf{TURMA} \\
            \hline
            #2 & \dataNun & \turma \\
            \hline
        \end{tabular}
    \end{center}
    
    #3
    
    \vspace{15pt}
}

\newcommand{\subtitulolis}{%
  \subtitulobase{ASSUNTO}{\assunto}{}%
}

\newcommand{\subtitulopr}{%
  \subtitulobase{ASSUNTO}{\assunto}{\boxed{\numProva}}%
}

\newcommand{\subtituloas}{%
  \subtitulobase{DISCIPLINA}{MATEMÁTICA}{\boxed{\numProva}}%
}

\newcommand{\subtituload}{%
  \subtitulobase{DISCIPLINA}{MATEMÁTICA}{}%
}

\newcommand{\valorquestao}{}

\newcommand{\definevalorquestao}{%
    \ifdefstring{\atividade}{lista}{%
        % Não faz nada para lista
    }{%
        \ifdefstring{\modo}{objetivo}{%
            \renewcommand{\valorquestao}{\ValorQObj ponto}
        }{%
            \renewcommand{\valorquestao}{\ValorQDisc pontos}
        }%
    }%
}

\newenvironment{exercicioBanco}[1][]{%
    \ifdefstring{\atividade}{lista}{%
        \begin{exercicio}%
    }{%
        \begin{exercicio}[#1]%
    }%
}{%
    \end{exercicio}
}

\ifdefstring{\atividade}{lista}{
    \pagecolor{background-color}
}{
    
}



\hypersetup{
    pdfauthor={\professor},
    pdftitle={\titulo \-- \assunto \-- \turma \-- \numProva},
}

\title{\titulo \-- \assunto \-- \turma \-- \numProva}
\author{\professor}
\date{\data}

%************* INÍCIO DO DOCUMENTO *************
\begin{document}
    \NoBgThispage
    \pagenumbering{alph}
    %
        
    \phantomsection
    \addcontentsline{toc}{section}{Capa}
    \input{../../AT_CapaAberta}
    \newpage
    \null
    \thispagestyle{empty}
    \addtocounter{page}{1}
    \newpage
        
    \NoBgThispage
    \pagenumbering{arabic}
    %
        
    \subtitulopr
    
    \vspace{15pt}
    \newcommand{\modo}{objetivo} % objetivo ou discursivo
    
    \phantomsection
    \addcontentsline{toc}{section}{Exercício 1}
    \input{../FQ_Exercicio2_1}
    \vspace{15pt}
    
    \phantomsection
    \addcontentsline{toc}{section}{Exercício 2}
    %%%%%%%%%%%%%%%%%%%%%%%%%%%%%%%%%
%
%       EXERCÍCIO
%
%%%%%%%%%%%%%%%%%%%%%%%%%%%%%%%%%

\ifdefstring{\atividade}{prova}{%
    \ifdefstring{\modo}{objetivo}{%
        \renewcommand{\valorquestao}{\ValorQObj\ ponto}
    }{%
        \renewcommand{\valorquestao}{\ValorQDisc\ pontos}
    }%
}%

\begin{exercicioBanco}[\valorquestao]
Uma das maneiras de saber se o seu peso está adequado à sua altura é calculando o Índice de Massa Corporal (IMC), que pode ser obtido através da fórmula dada por
%
\[
    h^{2}\cdot \text{IMC} - P = 0,
\]
%
sendo \(h\) a altura da pessoa em metros (m) e \(P\) o seu peso em quilogramas (kg). Se uma pessoa possui IMC igual a 40 kg/m\(^2\) e está com peso igual a 120 kg, indique a alternativa que apresenta o valor mais aproximado de sua altura em metros.

% Define as alternativas
\newcommand{\alternativas}{%
    \begin{center}
        \begin{tabularx}{\textwidth}{XXXXX}
            (a) \(1,4\). &
            (b) \(1,5\). &
            (c) \(1,6\). &
            (d) \(1,7\). &
            (e) \(1,8\).
        \end{tabularx}
    \end{center}
}

% Define a resposta correta
\newcommand{\resposta}{D}

% Lógica condicional para exibição
\ifdefstring{\atividade}{lista}{%
    \alternativas
    \vspace{0.5em}
    
    \noindent\textbf{Resposta:} letra \textbf{\resposta}.
}{%
    \ifdefstring{\modo}{objetivo}{%
        \alternativas
    }{%
        % modo = discursiva → não mostra alternativas
    }
}
\end{exercicioBanco}


    \vspace{15pt}
    
    \phantomsection
    \addcontentsline{toc}{section}{Exercício 3}
    %%%%%%%%%%%%%%%%%%%%%%%%%%%%%%%%%
%
%       EXERCÍCIO
%
%%%%%%%%%%%%%%%%%%%%%%%%%%%%%%%%%

\definevalorquestao

\begin{exercicioBanco}[\valorquestao]
Seja \(p(x) = mx^{2} + nx + 1\). Se \(p(2) = 0\) e \(p(-1) = 0\), então os valores de \(m\) e \(n\) são, respectivamente, iguais a?

% Define as alternativas
\newcommand{\alternativas}{%
    \begin{center}
        \begin{tabularx}{\textwidth}{XXXXX}
            (a) \(-\dfrac{1}{2}\) e \(\dfrac{1}{2}\). &
            (b) \(-1\) e \(1\). &
            (c) \(1\) e \(\dfrac{1}{2}\). &
            (d) \(-1\) e \(-\dfrac{1}{2}\). &
            (e) \(1\) e \(-\dfrac{1}{2}\).
        \end{tabularx}
    \end{center}
}

% Define a resposta correta
\newcommand{\resposta}{A}

% Lógica condicional para exibição
\ifdefstring{\atividade}{lista}{%
    \alternativas
    \vspace{0.5em}
    
    \noindent\textbf{Resposta:} letra \textbf{\resposta}.
}{%
    \ifdefstring{\modo}{objetivo}{%
        \alternativas
    }{%
        % modo = discursiva → não mostra alternativas
    }
}
\end{exercicioBanco}


    \vspace{15pt}
    
    \phantomsection
    \addcontentsline{toc}{section}{Exercício 4}
    \input{../FQ_Exercicio11_1}
    \vspace{15pt}
    
    \phantomsection
    \addcontentsline{toc}{section}{Exercício 5}
    %%%%%%%%%%%%%%%%%%%%%%%%%%%%%%%%%
%
%       EXERCÍCIO
%
%%%%%%%%%%%%%%%%%%%%%%%%%%%%%%%%%

\definevalorquestao

\begin{exercicioBanco}[\valorquestao]
No Laboratório de Química do CEFET, após várias medidas, um estudante concluiu que a concentração de certa substância em uma amostra variava em função do tempo, medido em horas, segundo a função quadrática \(f(t) = 5t - t^{2}\). Determine em que momento, após iniciadas as medidas, a concentração dessa substância foi máxima nessa amostra.

% Define as alternativas
\newcommand{\alternativas}{%
    \begin{center}
        \begin{tabularx}{\textwidth}{XXXXX}
            (a) \(\dfrac{25}{2}\). &
            (b) \(\dfrac{25}{4}\). &
            (c) \(\dfrac{15}{4}\). &
            (d) \(\dfrac{15}{2}\). &
            (e) \(\dfrac{20}{3}\).
        \end{tabularx}
    \end{center}
}

% Define a resposta correta
\newcommand{\resposta}{B}

% Lógica condicional para exibição
\ifdefstring{\atividade}{lista}{%
    \alternativas
    \vspace{0.5em}
    
    \noindent\textbf{Resposta:} letra \textbf{\resposta}.
}{%
    \ifdefstring{\modo}{objetivo}{%
        \alternativas
    }{%
        % modo = discursiva → não mostra alternativas
    }
}
\end{exercicioBanco}


    \vspace{15pt}
    
    \phantomsection
    \addcontentsline{toc}{section}{Exercício 6}
    %%%%%%%%%%%%%%%%%%%%%%%%%%%%%%%%%
%
%       EXERCÍCIO
%
%%%%%%%%%%%%%%%%%%%%%%%%%%%%%%%%%

\definevalorquestao

\begin{exercicioBanco}[\valorquestao]
O custo para produzir um produto é \(C(x) = 50 + 2x + 0,1x^{2}\), onde \(x\) é a quantidade produzida do produto. Cada unidade é vendida por R\$\(6,50\). Entre que valores deve variar \(x\) para não haver prejuízo?

% Define as alternativas
\newcommand{\alternativas}{%
    \begin{center}
        \begin{tabularx}{\textwidth}{XXX}
            (a) \(x \le 20\) ou \(x \ge 25\). &
            (b) \(10 \le x \le 50\). &
            (c) \(x \ge 20\). \\[5pt]
            (d) \(x \le 25\). &
            (e) \(20 \le x \le 25\).
        \end{tabularx}
    \end{center}
}

% Define a resposta correta
\newcommand{\resposta}{E}

% Lógica condicional para exibição
\ifdefstring{\atividade}{lista}{%
    \alternativas
    \vspace{0.5em}
    
    \noindent\textbf{Resposta:} letra \textbf{\resposta}.
}{%
    \ifdefstring{\modo}{objetivo}{%
        \alternativas
    }{%
        % modo = discursiva → não mostra alternativas
    }
}
\end{exercicioBanco}


    \vspace{15pt}
    
    %\newpage
    %\phantomsection
    %\addcontentsline{toc}{section}{Rascunho}
    %\input{../../Rascunho}
    
\end{document}
%*************** FINAL DO DOCUMENTO ***************
